\scanpage{6}
\uline{Corollaire} Si \inserted{de plus} $G$ est de t.f. sur $A$, alors $G$ admet une représentation linéaire fidèle.
\figure{\begin{tikzpicture}[x=1cm,y=1cm]\draw (0,0)--(-.03,.65);\end{tikzpicture}}

Il existe $N$ tel que a) $N$ engendre l'algèbre $B$, b) $\tilde N=N$, c) $\Delta N\subset N\otimes B$. \deleted{\unsure{le}} Alors \deleted{\unsure{le}} \unsure{tt} pt de $G$ à valeurs dans $A'/A$ qui opère \unsure{trivialt} sur $N'$, opère \unsure{trivialt} sur $B'$, \deleted{\unsure{et par suite}} donc est \deleted{\unsure{sur}} l'unité. Donc la repr. est fidèle.

\uline{Corollaire} Si $A$ est un corps $k$, alors, \unsure{tt} groupe affine \deleted{\unsure{de}} $G$ sur $k$ est lim. \deleted{\unsure{inductive}} projective d'un syst. projectif strict (\unsure{morph.} de transition des morphismes fid plats) de groupes affines \uline{de t.f.}.
\mspage{5}
